% Начинаем с новой страницы, оформляем как главу без номера
\chapter*{ПЕРЕЧЕНЬ УСЛОВНЫХ ОБОЗНАЧЕНИЙ, СИМВОЛОВ И ТЕРМИНОВ}
\addcontentsline{toc}{chapter}{ПЕРЕЧЕНЬ УСЛОВНЫХ ОБОЗНАЧЕНИЙ, СИМВОЛОВ И ТЕРМИНОВ}

В настоящей пояснительной записке применяют следующие термины с соответствующими определениями, а также сокращения и условные обозначения:

\vspace{0.5cm}

\begin{description}
    \item[БГУИР] --- Белорусский государственный университет информатики и радиоэлектроники;
    \item[ПО] --- программное обеспечение;
    \item[ПЗ] --- пояснительная записка;
    \item[ТЗ] --- техническое задание;
    \item[БД] --- база данных;
    \item[API] --- программный интерфейс приложения (Application Programming Interface);
    \item[MQTT] --- протокол обмена сообщениями между устройствами (Message Queuing Telemetry Transport);
    \item[$\Delta t$] --- погрешность измерения температуры, $^\circ$C;
    \item[$\eta$] --- коэффициент полезного действия;
    \item[IoT] --- интернет вещей (Internet of Things).
\end{description}

\vspace{0.5cm}

% Если разделяете на категории, можно использовать подразделы без нумерации:
\section*{Определения}

\noindent \textbf{Жизненный цикл ПО} --- период времени, начинающийся с момента принятия решения о необходимости создания программного продукта и заканчивающийся в момент его полного изъятия из эксплуатации.

\newpage